\begin{description}
    \item[Angular] Framework JavaScript développé par Google pour créer des applications web single-page (SPA). Utilisé pour développer l'interface utilisateur de l'application.
    
    \item[API] Application Programming Interface - Interface de programmation permettant la communication entre différents composants logiciels.
    
    \item[CDS] Centre de Services - équipe de 9 personnes de la DR Sillon Rhodanien spécialisée dans le développement d'outils numériques pour les métiers d'Enedis.
    
    \item[CheckMarx] Outil d'analyse statique de code utilisé chez Enedis pour détecter les vulnérabilités de sécurité dans le code source.
    
    \item[CI/CD] Continuous Integration/Continuous Deployment - Pratique DevOps d'automatisation des tests et déploiements pour améliorer la qualité et la rapidité des livraisons.

    \item[DIR2S] Système d'information RH national d'Enedis contenant les données des collaborateurs utilisé comme source pour l'automatisation des imports.
    
    \item[Downtime] Période d'indisponibilité d'un service ou application durant les opérations de maintenance ou de déploiement.
    
    \item[DR] Direction Régionale - Division territoriale d'Enedis.
    \item Il existe 25 DR réparties sur le territoire français.
    
    \item[DRHTS] Direction des Ressources Humaines et Transformation Santé-Sécurité - Direction nationale d'Enedis ayant porté le projet de refonte de l'outil de gestion des effectifs.
    
    \item[DSI] Direction des Systèmes d'Information - Direction nationale d'Enedis responsable de l'infrastructure informatique et des systèmes d'information.
    
    \item[EBIOS] Expression des Besoins et Identification des Objectifs de Sécurité - Méthode française d'analyse et de gestion des risques en sécurité informatique.
    
    \item[EFA] Effectif Fin d'Année - Indicateur RH correspondant au nombre de collaborateurs comptabilisés en fin d'année civile.
    
    \item[EMP] Effectif Moyen Payé - Nouvel indicateur RH majeur correspondant à la moyenne des effectifs payés sur une période donnée, introduit récemment chez Enedis.
    
    \item[GitLab] Plateforme de développement collaboratif utilisée pour le versioning du code, le suivi des issues et la gestion des sprints.
    
    \item[MC2] MyConsoleCloud - Plateforme cloud AWS d'Enedis utilisée pour le déploiement et l'hébergement des applications.
    
    \item[Migration] Processus de transformation ou mise à jour de la structure d'une base de données, généralement automatisé via des scripts.
    
    \item[Mono-repository] Structure de code où plusieurs projets liés sont stockés dans un seul repository Git, facilitant la gestion des dépendances et la synchronisation.
    
    \item[NestJS] Framework Node.js pour construire des applications backend efficaces et évolutives.
    \item Utilisé pour développer l'API de l'application.
    
    \item[Node.js] Environnement d'exécution JavaScript côté serveur permettant d'exécuter du code JavaScript en dehors d'un navigateur web.
    
    \item[NG-ZORRO] Bibliothèque de composants UI pour Angular basée sur Ant Design, utilisée pour l'interface utilisateur de l'application.
    
    \item[OAuth2] Protocole d'autorisation standard utilisé pour l'authentification SSO (Single Sign-On) dans l'écosystème Enedis.
    
    \item[ORM] Object-Relational Mapping - Technique de programmation permettant de faire correspondre les bases de données relationnelles avec les objets d'un langage de programmation.
    
    \item[Pipeline] Ensemble d'étapes automatisées pour le build, test et déploiement d'une application dans un processus CI/CD.
    
    \item[Planning Poker] Technique d'estimation agile utilisant des cartes numérotées pour évaluer la complexité et l'effort requis pour développer les user stories.
    
    \item[POC] Proof of Concept - Prototype ou démonstration permettant de valider la faisabilité d'un concept ou d'une technologie.
    
    \item[Prince2] Projects IN Controlled Environments - Méthodologie de gestion de projet structurée utilisée pour la gouvernance et le pilotage stratégique.
    
    \item[Prisma] ORM moderne pour Node.js et TypeScript utilisé pour la gestion de la base de données et les requêtes SQL.
    
    \item[Recette] Phase de tests et validation d'une fonctionnalité par les utilisateurs métier avant sa mise en production définitive.
    
    \item[RGPD] Règlement Général sur la Protection des Données - Réglementation européenne sur la protection des données personnelles, contrainte majeure du projet.
    
    \item[Rollback] Processus de retour à une version précédente d'une application en cas de problème détecté après un déploiement.
    
    \item[Sprint] Période de développement courte (généralement 1-2 semaines) dans la méthodologie Agile durant laquelle l'équipe se concentre sur un ensemble de fonctionnalités défini.

    \item[SSO] Single Sign-On - Méthode d'authentification permettant à un utilisateur d'accéder à plusieurs applications avec une seule connexion.
    
    \item[TMA] Tierce Maintenance Applicative - Service de maintenance et d'évolution d'applications informatiques assuré par une équipe externe au développement initial.
    
    \item[TurboRepo / Nx] Ensemble d'outils pour mono-repositories.
    
    \item[TypeScript] Langage de programmation développé par Microsoft, extension typée de JavaScript apportant une meilleure robustesse et maintenabilité du code.
    
    \item[User Story] Description courte et simple d'une fonctionnalité du point de vue de l'utilisateur final, utilisée dans les méthodologies agiles pour définir les besoins.
\end{description}